\documentclass[10pt,a4paper]{amsart}
\usepackage[margin=19mm]{geometry}
\usepackage{amsmath,amssymb,mathtools}
\usepackage[scr=rsfs,cal=euler]{mathalfa}
\usepackage{xfrac}
\usepackage[unicode,hidelinks,bookmarks=false]{hyperref}
\newcommand{\ie}{i.e.\ }
\newcommand{\cf}{cf.\ }
\newcommand{\resp}{respectively\ }
\newcommand{\Subsection}{\subsection}
\providecommand{\nobreakdash}{\nobreak\hbox{-}}
\setlength{\emergencystretch}{2em}
\multlinegap0pt
\usepackage{mathtools}
\usepackage{amssymb}
\usepackage[shortlabels]{enumitem}
\usepackage{xcolor}
\newtheorem{theorem}{Theorem}
\newtheorem{lemma}{Lemma}
\usepackage{cleveref}
\crefname{theorem}{Theorem}{Theorems}
\crefname{lemma}{Lemma}{Lemmas}
\newcommand{\SO}{\operatorname{SO}}
\newcommand{\acts}{\curvearrowright}
\title{On some aspects of discrete groups acting ergodically on the boundary}
\author{Subhadip Dey, Mikołaj Frączyk, Sebastian Hurtado}
\date{Source: arXiv:2608.27274v1 (CC BY 4.0)}
\begin{document}
\maketitle

\noindent\emph{Source and licence.} On some aspects of discrete groups acting ergodically on the boundary. Version arXiv:2608.27274v1 (CC BY 4.0), distributed by arXiv under the Creative Commons Attribution 4.0 International licence, http://creativecommons.org/licenses/by/4.0/, as recorded in the arXiv licence metadata for this exact version and shown on the abstract page https://arxiv.org/abs/2608.27274v1. Full licence notice: \texttt{licenses/arxiv-2608.27274-CC-BY-4.0.txt}.\par\smallskip

\noindent\textit{Editorial scope.} Counted unit: the article's theorem thm:proper-on-V (source0.tex lines 574--576). The article's complete original proof of it is delivered with it (source0.tex lines 645--696, one \texttt{proof} environment of 52 lines, which the article places away from the statement). No mathematics is written, repaired or re-derived: the statement and the proof are the article's own lines, in the article's own notation, and no retained line is substituted or re-flowed.  Retained uncounted as the source-local context the proof needs: lines 586--597; lines 605--613.  Nothing was omitted from any retained span, so the package carries no omission record.  No retained line contains a citation, so the wrapper carries no bibliography and no bibliography entry.  No retained line contains a figure environment, an image-inclusion command or drawing code, so no image is delivered and the package declares no asset dependency.  Editorial formatting only, in addition: an AMS \texttt{amsart} preamble that restates the article's own declarations and macros used by the retained lines, namely the environments \texttt{theorem}, \texttt{lemma} on separate counters, together with the standard packages those lines need.  The retained environments print in this document as Theorem 1, Lemma 1 in order of appearance, whereas the article prints the same items with the numbering of its own full document.  The complete native source of the article as distributed in the arXiv e-print for this exact version is bundled unchanged as \texttt{sources/208704/source0.tex}.
\begin{lemma}\label{lem:T}
    Let $(c_n)$ be a sequence in $G$ converging to $c\in G$, let $(t_n)$ be a sequence in $T$, and let $\delta_n \coloneqq c_n t_n$. Suppose that $\sigma_1(\delta_n)/\sigma_2(\delta_n) \to\infty$ as $n \to\infty$. Then, after passing to a subsequence,
    \[
     \frac{\delta_n}{\sigma_1(\delta_n)} \to B,
    \]
    where $B$ is a rank one matrix satisfying
    \[
     \operatorname{im} B = c(\mathbb R e_1)
     \quad \text{or} \quad
     \operatorname{im} B^\dagger = \mathbb R e_1.
    \]
\end{lemma}

    \begin{equation}\label{eqn1:lem:pd}
    t(a,u)=
    \begin{pmatrix}
    a&-a\,u^\flat&-\frac a2q_W(u)\\
    0&I&u\\
    0&0&a^{-1}
    \end{pmatrix},
    \qquad a>0,\ u\in W,
    \end{equation}

\begin{theorem}\label{thm:proper-on-V}
Let $\Gamma<G=\SO(n,2)$ be a  $Q$-Anosov subgroup. Then $\Gamma$ acts freely and properly discontinuously on $V_\Gamma$ (with the topology coming from $d_{G/T}$).
\end{theorem}

\begin{proof}[Proof of \Cref{thm:proper-on-V}]
    Recall that $V_\Gamma = \{ g T :\ g([e_1]) \not\in\Lambda_\Gamma \}$.

    We first show that the action $\Gamma \acts V_\Gamma$ is free. Suppose that $\gamma g T = g T$ for some $\Gamma \in\Gamma$. Then $t \coloneqq g^{-1} \gamma g \in T$. Writing $t = t(a,u)$, the upper triangular matrix given by \eqref{eqn1:lem:pd}, we have
    \[
     \operatorname{Spec} (t) = \{ a,\underbrace{1,\dots,1}_{n \text{ times}}, a^{-1} \}.
    \]

    If $\gamma$ is of finite order, the so does $t$. Hence, $a = a^{-1} = 1$ and so $t$ is unipotent. Note that $t^k = t(1, ku)$. Since $t$ is of finite order, $u=0$, and hence, $\gamma = e$.

    Else, $\gamma$ is of infinite order. Since $\langle \gamma\rangle$ is $Q$-Anosov, it is proximal in the standard representation; so, $a \ne 1$. If $a>1$, then $[e_1]$ is the attracting fixed point of $t$. Thus, $g[e_1]$ is the attracting fixed point of $\gamma$, which must lie in $\Lambda_\Gamma$, contradicting $gT \in V_\Gamma$. Similarly, if $a<1$, then $g[e_1]$ is the repelling fixed point of $\gamma$, which lies in $\Lambda_\Gamma$, again giving a contradiction.

    This concludes the proof of freeness of the action.

    We now show that the action $\Gamma \acts V_\Gamma$ is properly discontinuous.
    If not, there exist sequences $(\gamma_n)$ in $\Gamma$ and $(g_n)$ in $G$ and $g,h\in G$ such that $(\gamma_n)$ consists of pairwise distinct elements and
    \begin{equation}\label{eqn2:lem:pd}
     g_n T \to g T \in V_\Gamma
     \quad\text{and}\quad
     \gamma_n g_n T \to h T \in V_\Gamma
    \end{equation}
    as $n\to\infty$. For $n\in\mathbb N$, write $\gamma_n g_n = h_n t_n$, for some $h_n \in G$ and $t_n \in T$.
    Using local sections of $G \to G/T$, we can find appropriate representatives to further assume that $g_n \to g$ and $h_n \to h$ as $n\to\infty$.

    For $n\in\mathbb N$, let
    \[
     \delta_n \coloneqq g_n^{-1} \gamma_n g_n = c_n t_n,
    \]
    where $c_n = g_n^{-1} h_n$. Clearly, $c_n \to c \coloneqq g^{-1}h$ as $n\to\infty$. Since $(g_n)$ is bounded, the $Q$-Anosov property ensures that
    $\sigma_1(\delta_n)/\sigma_2(\delta_n) \to\infty$ as $n\to\infty$. 
    
    Thus, by \Cref{lem:T}, we can pass to a subsequence so that
    \[
     \frac{\delta_n}{\sigma_1(\delta_n)} \to B,
    \]
    for a rank one matrix $B$ satisfying either $\operatorname{im} B = c(\mathbb R e_1)$ or $\operatorname{im} B^\dagger = \mathbb R e_1$. 
    
    In the first case,
    \[
     \frac{\gamma_n}{\sigma_1(\delta_n)} = \frac{1}{\sigma_1(\delta_n)} g_n \delta_n g_n^{-1} \to g B g^{-1},
    \]
    as $n\to\infty$.
    Since $\operatorname{im} {g B g^{-1}} = gc (\mathbb R e_1) = h (\mathbb R e_1)$, we get that $h([e_1]) \in \Lambda_\Gamma$. This contradicts the assumption that $h T \in V_\Gamma$ (see \eqref{eqn2:lem:pd}).

    In the second case, we have that 
    \[
     \frac{\gamma_n^{-1}}{\sigma_1(\delta_n)} = \frac{\gamma_n^\dagger}{\sigma_1(\delta_n)} =  \frac{1}{\sigma_1(\delta_n)} g_n \delta_n^\dagger g_n^{-1} \to g B^\dagger g^{-1},
    \]
    as $n\to\infty$. This also gives, similar to the preceding paragraph, $g([e_1]) \in \Lambda_\Gamma$, again contradicting the assumption that $gT \in V_\Gamma$ (see \eqref{eqn2:lem:pd}).

    The proof of the theorem is now complete.
\end{proof}

\end{document}
